\chapter*{Perspectives}
\addcontentsline{toc}{chapter}{Perspectives}

Les évolutions suivantes prolongent le travail réalisé sans en changer l'architecture :
\begin{itemize}
    \item \textbf{Automatiser la vérification} : ajouter une suite de tests des services du processus
          principal (authentification, portée des secrétaires, calcul des moyennes, génération des
          rapports) et des tests de rendu des écrans critiques, puis l'exécuter avant chaque
          publication ;
    \item \textbf{Superviser le serveur depuis l'application} : exposer dans la page de diagnostic
          l'état du service PostgreSQL, l'espace disque du partage, l'âge de la dernière sauvegarde
          réussie et une alerte lorsque la sauvegarde a échoué ;
    \item \textbf{Compléter l'internationalisation} : faire passer les messages d'erreur du processus
          principal par des codes traduisibles, afin que l'anglais ne reste plus la seule langue des
          refus et des validations ;
    \item \textbf{Assister la planification} : signaler automatiquement les conflits de salle, de
          groupe et d'enseignant avant l'enregistrement d'une séance, et proposer les créneaux libres
          qui respectent le volume horaire restant ;
    \item \textbf{Étendre la consultation hors du poste} : publier en lecture seule, sur le réseau
          local, l'emploi du temps et les notes d'un étudiant, sans ouvrir l'application complète ;
    \item \textbf{Généraliser la mise à jour automatique} : appliquer aux postes Linux la vérification
          par empreinte et l'installation depuis le partage déjà disponibles sur Windows, afin que
          l'ensemble du parc se mette à jour de la même manière.
\end{itemize}
